\metatitle{Eulerian polynomials}

\metadescription{Eulerian polynomials and numbers: descents, recurrences,
generating functions, Worpitzky's identity, real roots, gamma-positivity,
Pólya frequency sequences, asymptotic normality, and polytope connections.}

\metakeywords{Eulerian polynomial, Eulerian number, permutation descent,
Worpitzky identity, real-rooted polynomial, gamma-positive polynomial,
Pólya frequency sequence, permutohedron, hypersimplex}


\section[eulerian]{Eulerian polynomials}

For a permutation $\sigma=\sigma_1\sigma_2\dotsb\sigma_n\in\symS_n$, a
\defin{descent} is an index $i\in[n-1]$ such that
$\sigma_i>\sigma_{i+1}$.  Write $\des(\sigma)$ for the number of descents.

\begin{definition}[Eulerian polynomials and numbers]
\label{eulerianPolynomial}\label{eulerianPolynomialDefinition}
For $n\geq1$, the \defin{Eulerian polynomial} is
\[
 A_n(t)\coloneqq\sum_{\sigma\in\symS_n}t^{\des(\sigma)}
 =\sum_{k=0}^{n-1}A(n,k)t^k,
\]
where the \defin{Eulerian number} $A(n,k)$ counts permutations in $\symS_n$
with $k$ descents.  We set $A_0(t)=1$ and take $A(n,k)=0$ outside
$0\leq k\leq n-1$.
\end{definition}

The first polynomials are
\[
\begin{array}{c|l}
n&A_n(t)\\
\hline
0&1\\
1&1\\
2&1+t\\
3&1+4t+t^2\\
4&1+11t+11t^2+t^3\\
5&1+26t+66t^2+26t^3+t^4.
\end{array}
\]
Their coefficients form the Eulerian triangle \oeis{A008292}; the version
with a row for $n=0$ and a terminal zero is \oeis{A173018}.  The same numbers
enumerate permutations by excedances.  Inversion also shows that descents of
$\sigma^{-1}$ have the same distribution, although
$\des(\sigma)$ and $\des(\sigma^{-1})$ need not agree for an individual
permutation.

The classical monograph of Foata and Schützenberger
\cite{FoataSchutzenberger1970} develops the combinatorial theory of Eulerian
polynomials.  Petersen's modern account \cite{Petersen2015EulerianNumbers}
continues through gamma-positivity, polytopes, and Coxeter groups.  See also
\cite[Sec.~1.4]{StanleyEC1}.


\subsection[eulerianRecurrences]{Recurrences}

Inserting $n$ into a permutation of $[n-1]$ gives
\begin{equation}
 A(n,k)=(k+1)A(n-1,k)+(n-k)A(n-1,k-1).
\end{equation}
Indeed, a permutation with $k$ descents has $k+1$ insertion slots that
preserve this number.  A permutation with $k-1$ descents has $n-k$ slots
that create one new descent.
Multiplying by $t^k$ and summing gives the differential recurrence
\begin{equation}
 A_{n+1}(t)=(1+nt)A_n(t)+t(1-t)A'_n(t),
 \qquad A_0(t)=1.
\end{equation}
One also has the binomial recurrence
\[
 A_n(t)=\sum_{j=0}^{n-1}\binom{n}{j}A_j(t)(t-1)^{n-1-j}
 \qquad(n\geq1).
\]


\subsection[eulerianGeneratingFunctions]{Generating functions and power sums}

The exponential generating function is
\begin{equation}
 \sum_{n\geq0}A_n(t)\frac{x^n}{n!}
 =\frac{t-1}{t-\exp((t-1)x)}.
\end{equation}
As an identity of formal power series in $t$,
\begin{equation}
 \frac{A_n(t)}{(1-t)^{n+1}}
 =\sum_{j\geq0}(j+1)^n t^j.
\end{equation}
For $|t|<1$ this is also an analytic identity.  Thus, for $n\geq1$,
\[
 \operatorname{Li}_{-n}(t)
 =\sum_{j\geq1}j^nt^j
 =\frac{tA_n(t)}{(1-t)^{n+1}}.
\]
This normalization explains why some references attach a factor $t$ to the
Eulerian polynomial: their numerator is $tA_n(t)$ rather than $A_n(t)$.


\subsection[eulerianWorpitzky]{Worpitzky's identity and explicit formulas}

The Eulerian numbers are the connection coefficients in
\defin{Worpitzky's identity}
\begin{equation}
 x^n=\sum_{k=0}^{n-1}A(n,k)\binom{x+k}{n}.
\end{equation}
Equivalently, inclusion--exclusion gives
\begin{equation}
 A(n,k)=\sum_{j=0}^{k}(-1)^j\binom{n+1}{j}(k+1-j)^n.
\end{equation}
If $S(n,k)$ denotes a Stirling number of the second kind, then
\begin{equation}
 A_n(t)=\sum_{k=1}^{n}k!S(n,k)t^{k-1}(1-t)^{n-k}
 \qquad(n\geq1).
\end{equation}
The exponent $k-1$ in this formula is another place where the shifted and
unshifted Eulerian conventions must not be mixed.


\subsection[eulerianSymmetry]{Symmetry and special values}

Complementing every value, or reversing the one-line notation, exchanges
descents and ascents.  Hence
\begin{equation}
 A(n,k)=A(n,n-1-k),
 \qquad
 t^{n-1}A_n(1/t)=A_n(t).
\end{equation}
In particular, $A_n(t)$ is palindromic of degree $n-1$.  Two useful
evaluations are
\[
 A_n(1)=n!,
 \qquad
 A_{2m}(-1)=0 \quad(m\geq1).
\]
For odd indices,
\[
 A_{2m+1}(-1)=(-1)^m E_{2m+1},
\]
where $E_n$ is the Euler zigzag number \oeis{A000111}.  Thus the nonzero
absolute values are the tangent numbers \oeis{A000182}; see
\cite[Ch.~5]{FoataSchutzenberger1970}.


\subsection[eulerianZeros]{Zeros, interlacing, and limit laws}

\begin{theorem}[Frobenius, \cite{Frobenius1910}]
For $n\geq2$, the polynomial $A_n(t)$ has $n-1$ simple, strictly negative
zeros.  Moreover, the zeros of $A_n(t)$ strictly interlace those of
$A_{n+1}(t)$.
\end{theorem}

The differential recurrence above gives a short inductive proof; see the
\hyperref[ex:eulerianSturm]{Eulerian Sturm-sequence example} for the root
argument.  Real-rootedness implies log-concavity and unimodality of every
Eulerian row.  By the
\hyperref[aissenSchoenbergWhitney]{Aissen--Schoenberg--Whitney theorem}, each
row $(A(n,0),\dotsc,A(n,n-1))$, extended by zeros, is also a
\hyperref[polyaFrequencySequenceDefinition]{Pólya frequency sequence}.
Consequently its Toeplitz matrix is totally non-negative; this is a TNN
statement, not a strict total-positivity statement.

If $D_n$ is the number of descents of a uniformly random permutation in
$\symS_n$, then, for $n\geq2$,
\[
 \operatorname{E}(D_n)=\frac{n-1}{2},
 \qquad
 \operatorname{Var}(D_n)=\frac{n+1}{12}.
\]
Furthermore,
\[
 \frac{D_n-(n-1)/2}{\sqrt{(n+1)/12}}
 \mathrel{\xrightarrow{\mathrm d}}\mathcal N(0,1).
\]
This is a classical instance of the
\hyperref[realRootedNormality]{central and local limit theory} for
real-rooted coefficient arrays; see \cite{Bender1973}.


\subsection[eulerianGamma]{Gamma-positivity}

The palindromic symmetry above can be strengthened.  For $n\geq1$ there are
non-negative integers $\gamma_{n,k}$ such that
\begin{equation}
 A_n(t)=\sum_{k=0}^{\lfloor(n-1)/2\rfloor}
 \gamma_{n,k}t^k(1+t)^{n-1-2k}.
\end{equation}
Thus $A_n(t)$ is \hyperref[gammaPositivePolynomial]{gamma-positive}; its
gamma-vectors form \oeis{A101280}.  The
\hyperref[foataStrehlAction]{Foata--Strehl valley-hopping action} gives a
combinatorial proof and interpretation of the coefficients
\cite{FoataSchutzenberger1970}.  See
\cite[Ch.~4]{Petersen2015EulerianNumbers} for a modern account.


\subsection[eulerianGeometry]{Polytope connections}

Eulerian polynomials occur naturally as $h$- and $h^*$-polynomials.
\begin{itemize}
\item The $h$-polynomial of the
\hyperref[permutohedron]{permutohedron $\operatorname{Perm}_n$} is $A_n(t)$.
Equivalently, the Eulerian numbers form its $h$-vector.

\item The \hyperref[hStarPolynomial]{$h^*$-polynomial} of the cube
$[0,1]^n$ is $A_n(t)$.  Indeed, its Ehrhart series is
\[
 \sum_{m\geq0}(m+1)^nt^m=\frac{A_n(t)}{(1-t)^{n+1}},
\]
which is the power-sum identity above.

\item The normalized volume of the
\hyperref[hypersimplex]{hypersimplex $\Delta_{k,n}$} is
$A(n-1,k-1)$.  More refined descent--excedance formulas describe its
$h^*$-polynomial; see \cite{Li2012}.
\end{itemize}


\subsection[eulerianGeneralizations]{Generalizations}

Several standard extensions retain part of this structure.  Joint
distributions such as
\[
 \sum_{\sigma\in\symS_n}q^{\maj(\sigma)}t^{\des(\sigma)}
\]
give $q$-Eulerian, or Euler--Mahonian, polynomials and specialize to $A_n(t)$
at $q=1$; conventions vary, so the statistic and shift should always be
stated.  See \cite{ShareshianWachs2007} and the
\hyperref[eulerianSymmetric]{Eulerian quasisymmetric-function page}.

For a finite Coxeter group $W$, one similarly considers
$\sum_{w\in W}t^{\des_W(w)}$.  The classical polynomials above are type
$A$; the type $B$ and type $D$ coefficient triangles are respectively
\oeis{A060187} and \oeis{A066094}.  See the
\hyperref[ex:typeBEulerianSturm]{type $B$ and $D$ real-rootedness discussion}
and the \hyperref[coxeterTypeB]{Coxeter-groups page}.  Descent polynomials of
linear extensions give the broader class of
\hyperref[pEulerianPolynomial]{$P$-Eulerian polynomials}.
